\section{Introduction}

A proper coloring of a graph assigns different colors to adjacent vertices.
An independent set is a set of pairwise nonadjacent vertices.
For a nonempty finite graph $G$, write $n=|V(G)|$ and let $\alpha(G)$
be its maximum independent-set size.
Every three-colorable graph contains an independent set of size at least
$n/3$: one of its three color classes has that size.
We prove that distinguishing this structured case from graphs with an
arbitrarily small fixed independence ratio is NP-hard.

\begin{theorem}\label{thm:main}
For every fixed real $0<\delta<1/3$, there is a deterministic
polynomial-time algorithm $R_\delta$ that maps each explicitly encoded
Boolean $3$-CNF formula $\phi$ to a nonempty finite simple undirected
unweighted graph $G_\phi$ such that
\[
 \begin{array}{ll}
 \phi\text{ is satisfiable}
   &\Longrightarrow G_\phi\text{ is three-colorable},\\[2pt]
 \phi\text{ is unsatisfiable}
   &\Longrightarrow \alpha(G_\phi)<\delta\,|V(G_\phi)|.
 \end{array}
\]
The running time and explicit output size are polynomial in the ordinary
binary encoding length of $\phi$. Their constants and polynomial degree
may depend on the fixed $\delta$, but not on $\phi$.
The coloring in the first implication covers every vertex and is not
supplied with the graph.
\end{theorem}

It suffices to prove Theorem~\ref{thm:main} for rational $\delta$.
For a fixed real threshold, choose a fixed rational $0<\delta_0<\delta$
and use $R_{\delta_0}$. We assume $\delta$ rational in the construction
and its analysis.
The density conclusion is stronger than an arbitrarily large fixed
lower bound on chromatic number. Indeed, small independence ratio forces
large chromatic number, whereas adjoining isolated vertices preserves
chromatic number and can make the independence ratio arbitrarily close
to one.

Approximate graph coloring asks, for fixed integers $3\le k\le c$,
to find a proper $c$-coloring of a graph promised to be $k$-colorable.
Its decision version distinguishes $k$-colorable graphs from graphs
that are not $c$-colorable. The following consequence recovers the
constant-palette hardness theorem.

\begin{corollary}\label{cor:coloring}
For every fixed pair of integers $3\le k\le c$, it is NP-hard to
distinguish $k$-colorable graphs from graphs that are not $c$-colorable.
\end{corollary}

\begin{proof}
Choose a rational $0<\delta<\min\{1/3,1/c\}$ and apply
Theorem~\ref{thm:main}. A three-colorable graph is $k$-colorable.
A $c$-colorable graph has an independent set of size at least $n/c$,
which the soundness bound excludes.
\end{proof}

\subsection{History and comparison with prior work}

The difficulty of approximate coloring persists even with a very small
promised palette. Khanna, Linial, and Safra established hardness of
four-coloring three-colorable graphs \cite{KhannaLinialSafra2000};
Guruswami and Khanna later gave a different proof and obtained
bounded-degree and edge-error variants \cite{GuruswamiKhanna2004}.
The algebraic theory of promise constraint satisfaction developed by
Barto, Bul\'{\i}n, Krokhin, and Opr\v{s}al gave hardness of
$(2k-1)$-coloring $k$-colorable graphs, and hence five colors when
$k=3$ \cite[Theorem~6.5]{BartoBulinKrokhinOprsal2021}.
Their Example~2.9 records the conjectured hardness for every fixed
$3\le k\le c$.

Guruswami and Sandeep showed that the ordinary $d$-to-$1$ Games
Conjecture with perfect completeness, for any fixed $d\ge2$, implies
constant-palette hardness for three-colorable graphs
\cite[Theorem~1]{GuruswamiSandeep2020}.
Their reduction first gives arbitrarily small independent-set density
with a $2d$-colorable completeness promise
\cite[Theorem~9]{GuruswamiSandeep2020}. The arc-graph reduction of
Krokhin, Opr\v{s}al, Wrochna, and \v{Z}ivn\'y
\cite[Theorem~1.8]{KrokhinOprsalWrochnaZivny2023} then reduces the
promised palette to three for the coloring conclusion. The resulting
three-colorable hardness statement does not include the independent-set
guarantee.
The companion article~\cite[Theorem~1.1]{OpenAIPerfectCompleteness2026}
establishes the ordinary perfect-completeness premise for $d=2$.

Fei, Minzer, and Wang have proved the $4$-to-$1$ Games Conjecture
with perfect completeness \cite[Theorem~1.6]{FeiMinzerWang2026}.
Combining it with the preceding reductions resolves the constant-palette
conjecture: their Corollary~1.7 gives the three-colorable versus not
$c$-colorable gap for every fixed $c$. Their Corollary~1.9 gives
arbitrarily small independent-set density with an eight-colorable
completeness promise, corresponding to $2d=8$.
Theorem~\ref{thm:main} combines this stronger form of soundness with
three-colorability of the entire graph. Our proof starts with ordinary
projection Label Cover and does not require a bound on projection-fiber
sizes.

The distinction between colorability and independence density also
appears in earlier conditional results. Dinur, Mossel, and Regev obtained
three-colorable completeness and arbitrarily small independent-set
soundness from their fish-shaped Label Cover conjecture
\cite[Conjecture~4.8 and Sections~4.3--4.5]{DinurMosselRegev2009}.
Braverman, Khot, Lifshitz, and Minzer obtained the same graph promise
from the Rich $2$-to-$1$ Games Conjecture using their invariance principle
for the multi-slice \cite[Corollary~1.19]{BravermanKhotLifshitzMinzer2022}.
The Rich condition requires that, at each left question, a uniformly
random incident constraint induce a uniformly random partition of the
left alphabet among all partitions into pairs. This is not part of the
ordinary $2$-to-$1$ premise above.

A different line of work relaxes completeness by permitting deletion
of a small fraction of vertices. Dinur, Khot, Perkins, and Safra proved
that, for arbitrarily small fixed $\varepsilon>0$, it is NP-hard to find
an independent set of density $1/9$ when an induced
three-colorable subgraph occupies at least $(1-\varepsilon)n$ vertices
\cite{DinurKhotPerkinsSafra2010}.
Hecht, Minzer, and Safra subsequently proved, under randomized reductions,
hardness of almost coloring such almost-three-colorable graphs with any
fixed palette \cite{HechtMinzerSafra2023}.
The all-vertex completeness in Theorem~\ref{thm:main} retains the original
coloring promise, while allowing every fixed positive soundness density.

Algorithmic work gives a complementary perspective on the remaining
quantitative gap. Bansal, Huang, and Lee give a polynomial-time algorithm
using $O(n^{0.19539})$ colors \cite{BansalHuangLee2026}, improving the
$\widetilde O(n^{0.19747})$ bound of Kawarabayashi, Thorup, and Yoneda
\cite{KawarabayashiThorupYoneda2024}.
In preprints posted after the September 24 version of this manuscript,
Narang and Tang report a randomized polynomial-time bound of
$O(n^{(13-\sqrt{97})/18+\varepsilon})$ for each fixed
$\varepsilon>0$ \cite{NarangTang2026}, and Anand reports a randomized
polynomial-time bound of $O(n^{4/23})$ \cite{Anand2026}.
Our reduction has fixed $\delta$ and fixed palettes; its polynomial
exponent can depend on these constants.

\subsection{Proof overview and method ancestry}

The starting problem is Label Cover: a bipartite collection of questions
and tests, each prescribing a projection from a left answer to a right
answer. The PCP theorem and parallel repetition supply instances in
which either every test can be satisfied or every labeling succeeds
on at most an arbitrarily small fixed fraction of tests
\cite{AroraEtAl1998,Raz1998,DinurSteurer2014}.

We replace one left question at a time by a right question, obtaining
chains of question tuples. This use of layered tuples and projection
constraints is related to the multilayered PCP construction of
Dinur, Guruswami, Khot, and Regev \cite{DinurGuruswamiKhotRegev2005}.
At each tuple, a point assigns a phase in $\T=\R/\mathbb Z$ to every
possible product answer. An answer projection $\pi:M\to M' $ pulls a
phase vector $x\in\T^{M'}$ back to $x\circ\pi\in\T^M$.
Link these two points with length zero, and link points within each
phase space with their sup circle distance. Shortest paths through
these links define a pseudometric. Edges compare one point with
the coordinatewise half-shift of another. A satisfying Label Cover
labeling evaluates each phase assignment at one product answer.
This evaluation is nonexpanding, and the edge threshold forces adjacent
vertices to have circle phases more than $1/3$ apart. Three equal
half-open intervals of the circle therefore color all vertices properly.

For soundness, an independent set gives bounded functions on the tuple
phase spaces that change sign under a half-shift and satisfy a common
Lipschitz bound. Austin's dimension-independent approximation theorem
\cite{Austin2016}, a continuous counterpart of Friedgut's junta theorem
\cite{Friedgut1998}, approximates these functions using boundedly many
answer coordinates. The difficulty is that approximation under product
measure need not survive identifications of coordinates by projections.
We therefore choose coordinate lists for whole small subchains and
introduce independent rotation variables at every layer of each
subchain. All required compatibility properties are proved within this
paper. Projecting the selected coordinates to the last layer places
all lists in a common answer set. A shared answer for two separated
subchains supplies compatible answers to an intervening Label Cover
test, so low source value makes these intersections rare.

Two further steps turn this structural information into a density
bound. First, a distributional alignment lemma chooses a law on sets
of layers before seeing the lists. On most selected sets, each listed
label can be assigned to an index shared by all its occurrences.
The bound is independent of the label universe. The proof combines
finite minimax, Ramsey homogenization of finite probability patterns,
and an order-invariant random-array argument. The homogenization step
has antecedents in the Ramsey theory of random variables
\cite{TrotterWinkler1998}; we give the particular alignment argument
in full.

Second, we sample coefficients at finitely many levels between $0$ and
$1$, with geometrically decreasing probabilities. Enough layers make
a coefficient equal to $1$ likely. That layer absorbs auxiliary uniform
phase shifts without changing their distribution. Alignment makes
dependence on a zero-coefficient layer small; decreasing a positive
coefficient by one level bounds its probability of substantial
dependence. A dimension-independent bound on influential coordinates
then gives a small mean square, and hence a small independent-set density.

Only fixed finite rational data enter the graph construction.
The probability experiment is expanded exactly into unweighted vertices,
so the reduction is deterministic.
Section~\ref{sec:preliminaries} states the standard inputs and derives
the analytic coordinate bounds.
Section~\ref{sec:construction} constructs the graph with finite parameters
and proves completeness. Section~\ref{sec:lists} extracts the coordinate
lists from an independent set and decodes intersections of separated lists.
Section~\ref{sec:alignment} states distributional alignment and applies
it to these lists. Section~\ref{sec:coefficient} proves the coefficient
estimate under explicit inequalities, and Section~\ref{sec:parameters}
chooses all constants in dependency order and completes
Theorem~\ref{thm:main}. The self-contained proof of distributional
alignment is given in Section~\ref{sec:alignment-proof}.
